\title{Spectral Normalization for Generative Adversarial Networks}

\begin{abstract}
One of the challenges in the study of generative adversarial networks is the instability of its training. In this paper, we propose a novel weight normalization technique called spectral normalization to stabilize the training of the discriminator. Our new normalization technique is computationally light and easy to incorporate into existing implementations. We tested the efficacy of spectral normalization on CIFAR10, STL-10, and ILSVRC2012 dataset, and we experimentally confirmed that spectrally normalized GANs (SN-GANs) is capable of generating images of better or equal quality relative to the previous training stabilization techniques. The code with Chainer~\cite[]{tokui2015chainer}, generated images and pretrained models are available at \url{https://github.com/pfnet-research/sngan_projection}.
\end{abstract}

\section{Introduction}
\textit{Generative adversarial networks}~(GANs)~\cite[]{goodfellow2014generative} have been enjoying considerable success as a framework of generative models in recent years, and it has been applied to numerous types of tasks and datasets~\cite[]{radford2015unsupervised, salimans2016improved, ho2016generative, li2017adversarial}. In a nutshell, GANs are a framework to produce a model distribution that mimics a given target distribution, and it consists of a generator that produces the model distribution and a discriminator that distinguishes the model distribution from the target. The concept is to consecutively train the model distribution and the discriminator in turn, with the goal of reducing the difference between the model distribution and the target distribution measured by the best discriminator possible at each step of the training. GANs have been drawing attention in the machine learning community not only for its ability to learn highly structured probability distribution but also for its theoretically interesting aspects. For example,~\cite[]{nowozin2016f, uehara2016generative, mohamed2016learning} revealed that the training of the discriminator amounts to the training of a good estimator for the density ratio between the model distribution and the target. This is a perspective that opens the door to the methods of
\textit{implicit models}~\cite[]{mohamed2016learning, tran2017deep} that can be used to carry out variational optimization without the direct knowledge of the density function.

A persisting challenge in the training of GANs is the performance control of the discriminator. In high dimensional spaces, the density ratio estimation by the discriminator is often inaccurate and unstable during the training, and generator networks fail to learn the multimodal structure of the target distribution. Even worse, when the support of the model distribution and the support of the target distribution are disjoint, there exists a discriminator that can perfectly distinguish the model distribution from the target~\cite[]{arjovsky2017towards}.
Once such discriminator is produced in this situation, the training of the generator comes to complete stop, because the derivative of the so-produced discriminator with respect to the input turns out to be 0. This motivates us to introduce some form of restriction to the choice of the discriminator.

In this paper, we propose a novel weight normalization method called \textit{spectral normalization} that can stabilize the training of discriminator networks. Our normalization enjoys following favorable properties.
\begin{tight_itemize}
\item Lipschitz constant is the only hyper-parameter to be tuned, and the algorithm does not require intensive tuning of the only hyper-parameter for satisfactory performance.
\item Implementation is simple and the additional computational cost is small.
\end{tight_itemize}
In fact, our normalization method also functioned well even without tuning Lipschitz constant, which is the only hyper parameter. In this study, we provide explanations of the effectiveness of spectral normalization for GANs against other regularization techniques, such as weight normalization~\cite[]{salimans2016weight}, weight clipping~\cite[]{arjovsky2017wasserstein}, and gradient penalty~\cite[]{gulrajani2017improved}.
We also show that, in the absence of complimentary regularization techniques (e.g., batch normalization, weight decay and feature matching on the discriminator), spectral normalization can improve the sheer quality of the generated images better than weight normalization and gradient penalty.

\section{Method}
In this section, we will lay the theoretical groundwork for our proposed method. Let us consider a simple discriminator made of a neural network of the following form, with the input ${\bm x}$:

\begin{align}
	f({\bm x}, \theta) = W^{L+1} a_L (W^L (a_{L-1}(W^{L-1}(\dots a_1(W^1{\bm x})\dots)))),
\end{align}

where $\theta:=\{W^1,\dots, W^L, W^{L+1} \}$ is the learning parameters set, $W^l\in \mathbb{R}^{d_l \times d_{l-1}}$, $W^{L+1}\in \mathbb{R}^{1\times{d_L}}$, and $a_l$ is an element-wise non-linear activation function. We omit the bias term of each layer for simplicity. The final output of the discriminator is given by

\begin{align}
	D({\bm x}, \theta) =  {\rm \mathcal{A}}(f({\bm x}, \theta)),
\end{align}

where $\mathcal{A}$ is an activation function corresponding to the divergence of distance measure of the user's choice. The standard formulation of GANs is given by
$$\min_{G} \max_{D} V(G, D)$$ where min and max of $G$ and $D$ are taken over the set of generator and discriminator functions, respectively. The conventional form of $V(G,D)$~\cite[]{goodfellow2014generative} is given by
$\mathrm{E}_{{\bm x}\sim q_{\rm data}} [\log D({\bm x})] + \mathrm{E}_{{\bm x}'\sim p_G} [\log(1-D({\bm x}'))], \label{eq:stdgan}$
where $q_{\rm data}$ is the data distribution and $p_G$ is the (model) generator distribution to be learned through the adversarial min-max optimization. The activation function ${\rm \mathcal{A}}$ that is used in the $D$ of this expression is some continuous function with range $[0,1]$ (e.g, sigmoid function). It is known that, for a fixed generator $G$, the optimal discriminator for this form of $V(G,D)$ is given by $D^*_G({\bm x}) := q_{\rm data}({\bm x})/(q_{\rm data}({\bm x}) + p_{G}({\bm x}))$.

The machine learning community has been pointing out recently that the function space from which the discriminators are selected crucially affects the performance of GANs. A number of works~\cite[]{uehara2016generative, qi2017loss, gulrajani2017improved} advocate the importance of Lipschitz continuity in assuring the boundedness of statistics. For example, the optimal discriminator of GANs on the above standard formulation takes the form

\begin{align}
	D_{G}^*({\bm x}) & = \frac{q_{\rm data}({\bm x})}{q_{\rm data}({\bm x})+p_G({\bm x})} = {\rm sigmoid}(f^*({\bm x})),
	\text{where }f^*({\bm x}) = \log q_{\rm data}({\bm x}) - \log p_G({\bm x}),
\end{align}

and its derivative

\begin{align}
\nabla_{\bm x} f^*(x) = \frac{1}{q_{\rm data}({\bm x})}\nabla_{{\bm x}}q_{\rm data}({\bm x}) - \frac{1}{p_G({\bm x})}\nabla_{{\bm x}}p_G({\bm x})
\end{align}

can be unbounded or even incomputable. This prompts us to introduce some regularity condition to the derivative of $f({\bm x})$.

A particularly successful works in this array are~\cite[]{qi2017loss, arjovsky2017wasserstein, gulrajani2017improved}, which proposed methods to control the Lipschitz constant of the  discriminator by adding regularization terms defined on input examples ${\bm x}$. We would follow their footsteps and search for the discriminator $D$ from the set of $K$-Lipschitz continuous functions, that is,

\begin{align}
	\mathop{\rm arg~max}\limits_{\|f\|_{\rm Lip}\leq K} V(G, D), \label{eq:lipschitz}
\end{align}

where we mean by $\|f\|_{\rm Lip}$ the smallest value $M$ such that $\|f({\bm x})-f({\bm x}')\| / \|{\bm x}-{\bm x}'\|\leq M$ for any ${\bm x},{\bm x}'$, with the norm being the $\ell_2$ norm.

While input based regularizations allow for relatively easy formulations based on samples, they also suffer from the fact that, they cannot impose regularization on the space outside of the supports of the generator and data distributions without introducing somewhat heuristic means. A method we would introduce in this paper, called \textit{spectral normalization}, is a method that aims to skirt this issue by \textit{normalizing} the weight matrices using the technique devised by~\cite{yoshida2017spectral}.

\subsection{\label{subsec:SN}Spectral Normalization} Our spectral normalization controls the Lipschitz constant of the discriminator function $f$ by literally constraining the spectral norm of each layer $g: {\bm h}_{in} \mapsto {\bm h}_{out}$. By definition, Lipschitz norm $\|g\|_{\rm Lip}$ is equal to $\sup_{\bm h} \sigma(\nabla g({\bm h}))$, where $\sigma(A)$ is the spectral norm of the matrix $A$ ($L_2$ matrix norm of $A$)

\begin{align}
	\sigma(A) := \max_{{\bm h}: {\bm h} \neq {\bm 0}} \frac{\|A{\bm h}\|_2}{\|{\bm h}\|_2} = \max_{\|{\bm h}\|_2\leq 1} \|A{\bm h}\|_2, \label{eq:spnorm}
\end{align}

which is equivalent to the largest singular value of $A$. Therefore, for a linear layer $g({\bm h}) = W{\bm h}$, the norm is given by $\|g\|_{\rm Lip} = \sup_{\bm h} \sigma(\nabla g({\bm h})) = \sup_{\bm h} \sigma(W) = \sigma(W) $. If the Lipschitz norm of the activation function $\|a_l\|_{\rm Lip}$ is equal to 1 \footnote{For examples, ReLU~\cite[]{jarrett2009best,nair2010rectified, glorot2011deep} and leaky ReLU~\cite[]{maas2013rectifier} satisfies the condition, and many popular activation functions satisfy $K$-Lipschitz constraint for some predefined $K$ as well.}, we can use the inequality $\|g_1\circ g_2\|_{\rm Lip}\leq\|g_1\|_{\rm Lip}\cdot\|g_2\|_{\rm Lip}$ to observe the following bound on $\|f\|_{\rm Lip}$:

\begin{align}
 	\|f\|_{\rm Lip} \leq & \|({\bm h}_{L}\mapsto W^{L+1}{\bm h}_{L})\|_{\rm Lip}\cdot \|a_L\|_{\rm Lip}\cdot\|({\bm h}_{L-1}\mapsto W^ L{\bm h}_{L-1})\|_{\rm Lip} \nonumber \\
    &\cdots \|a_1\|_{\rm Lip}\cdot\|({\bm h}_0\mapsto W^ 1{\bm h}_0)\|_{\rm Lip} = \prod_{l=1}^{L+1} \|({\bm h}_{l-1}\mapsto W^l{\bm h}_{l-1})\|_{\rm Lip} = \prod_{l=1}^{L+1} \sigma(W^l). \label{eq:ineq_lip}
\end{align}

Our \textit{spectral normalization} normalizes the spectral norm of the weight matrix $W$ so that it satisfies the Lipschitz constraint $\sigma(W) = 1$:

\begin{align}
\bar{W}_{\rm SN}(W) := W / \sigma(W) \label{eq:sn}.
\end{align}

If we normalize each $W^l$ using~\eqref{eq:sn}, we can appeal to the inequality~\eqref{eq:ineq_lip} and the fact that $\sigma\left(\bar{W}_{\rm SN}(W)\right) = 1$ to see that $\|f\|_{\rm Lip}$ is bounded from above by $1$.

Here, we would like to emphasize the difference between our spectral normalization and spectral norm "regularization" introduced by ~\cite{yoshida2017spectral}. Unlike our method, spectral norm "regularization" penalizes the spectral norm by adding explicit regularization term to the objective function. Their method is fundamentally different from our method in that they do not make an attempt to `set' the spectral norm to a designated value. Moreover, when we reorganize the derivative of our normalized cost function and rewrite our objective function \eqref{eq:grad_sn}, we see that our method is augmenting the cost function with a sample data \textit{dependent} regularization function. Spectral norm regularization, on the other hand, imposes sample data \textit{independent} regularization on the cost function,  just like L2 regularization and Lasso. 

\subsection{Fast Approximation of the Spectral Norm $\sigma(W)$}
As we mentioned above, the spectral norm $\sigma(W)$ that we use to regularize each layer of the discriminator is the largest singular value of $W$. If we naively apply singular value decomposition to compute the $\sigma(W)$ at each round of the algorithm, the algorithm can become computationally heavy. Instead, we can use the power iteration method to estimate $\sigma(W)$~\cite[]{golub2000eigenvalue, yoshida2017spectral}. With power iteration method, we can estimate the spectral norm with very small additional computational time relative to the full computational cost of the vanilla GANs. Please see Appendix~\ref{apdsec:algo_snorm} for the detail method and Algorithm ~\ref{algo:snorm} for the summary of the actual spectral normalization algorithm.

\subsection{Gradient analysis of the spectrally normalized weights \label{subsec:grad}}
The gradient\footnote{Indeed, when the spectrum has multiplicities, we would be looking at subgradients here. However, the probability of this happening is zero (almost surely), so we would continue discussions without giving considerations to such events.} of $\bar{W}_{\rm SN}(W)$ with respect to $W_{ij}$ is:

\begin{align}
	\frac{\partial \bar{W}_{\rm SN}(W)}{\partial W_{ij}} = \frac{1}{\sigma(W)}E_{ij} - \frac{1}{\sigma(W)^2} \frac{\partial \sigma(W)}{\partial W_{ij}} W &=\frac{1}{\sigma(W)}E_{ij} - \frac{[{\bm u}_{1}{\bm v}_{1}^{\rm T}]_{ij}} {\sigma(W)^2} W \\
    &= \frac{1}{\sigma(W)}\left(E_{ij} - [{\bm u}_{1}{\bm v}_{1}^{\rm T}]_{ij} \bar{W}_{\rm SN}\right),
\end{align}

where $E_{ij}$ is the matrix whose $(i,j)$-th entry is $1$ and zero everywhere else, and ${\bm u}_1$ and ${\bm v}_1$ are respectively the first left and right singular vectors of $W$. If ${\bm h}$ is the hidden layer in the network to be transformed by $\bar{W}_{SN}$, the derivative of the $V(G,D)$ calculated over the mini-batch with respect to $W$ of the discriminator $D$ is given by:

\begin{align}
 \frac{\partial V(G,D)}{\partial W} &= \frac{1}{\sigma(W)} \left(\hat{\mathop{\mathrm{E}}}\left[{\bm \delta} {\bm h}^{\rm T}\right] - \left(\hat{\mathop{\mathrm{E}}}\left[{\bm \delta}^{\rm T} \bar{W}_{\rm SN}{\bm h} \right] \right) {\bm u}_1 {\bm v}_1^{\rm T} \right) \\
 &=\frac{1}{\sigma(W)} \left(\hat{\mathop{\mathrm{E}}}\left[{\bm \delta} {\bm h}^{\rm T}\right] - \lambda {\bm u}_1 {\bm v}_1^{\rm T} \right) \label{eq:grad_sn}
\end{align}

where ${\bm \delta} := \left(\partial V(G,D) / \partial \left(\bar{W}_{\rm SN} {\bm h}\right)\right)^{\rm T}$, $\lambda:=\hat{\mathop{\mathrm{E}}}\left[{\bm \delta}^{\rm T} \left( \bar{W}_{\rm SN}{\bm h} \right)\right]$, and $\hat{\mathop{\mathrm{E}}}[\cdot]$ represents empirical expectation over the mini-batch.
$\frac{\partial V}{\partial W}=0$ when $\hat{\mathop{\mathrm{E}}}[{\bm \delta} {\bm h}^{\rm T}] = k{\bm u}_1 {\bm v}_1^T \label{eq:fixedpoint}$ for some $k\in \mathbb{R}$.

We would like to comment on the implication of~\eqref{eq:grad_sn}.
The first term $\hat{\mathop{\mathrm{E}}}\left[{\bm \delta} {\bm h}^{\rm T}\right]$ is the same as the derivative of the weights without normalization. In this light, the second term in the expression can be seen as the regularization term penalizing the first singular components with the \textit{adaptive} regularization coefficient $\lambda$.
$\lambda$ is positive when ${\bm \delta}$ and $\bar{W}_{\rm SN} {\bm h}$ are pointing in similar direction, and this prevents the column space of $W$ from concentrating into one particular direction in the course of the training. In other words, spectral normalization prevents the transformation of each layer from becoming to sensitive in one direction. We can also use spectral normalization to devise a new parametrization for the model. Namely, we can split the layer map into two separate trainable components: spectrally normalized map and the spectral norm constant. As it turns out, this parametrization has its merit on its own and promotes the performance of GANs (See Appendix~\ref{subsec:reparametrization}).

\section{Spectral Normalization vs Other Regularization Techniques}\label{sec:snvsothers}
The weight normalization introduced by~\cite{salimans2016weight} is a method that normalizes the $\ell_2$ norm of each row vector in the weight matrix. Mathematically, this is equivalent to requiring the weight by the weight normalization $\bar{W}_{\rm WN}$:

\begin{align}
	\sigma_1(\bar{W}_{\rm WN})^2 + \sigma_2(\bar{W}_{\rm WN})^2 + \cdots + \sigma_T(\bar{W}_{\rm WN})^2 = d_o,~{\rm where}~T=\min(d_i,d_o), \label{eq:svconstraint}
\end{align}

where $\sigma_t(A)$ is a $t$-th singular value of matrix $A$. Therefore, up to a scaler, this is same as the Frobenius normalization, which requires the sum of the squared singular values to be 1. These normalizations, however, inadvertently impose much stronger constraint on the matrix than intended. If $\bar{W}_{\rm WN}$ is the weight normalized matrix of dimension $d_i \times d_o$, the norm $\|\bar{W}_{\rm WN}{\bm h}\|_2$ for a fixed unit vector ${\bm h}$ is maximized at $\| \bar{W}_{\rm WN} {\bm h} \|_2 =\sqrt{d_o}$ when $\sigma_1(\bar{W}_{\rm WN})=\sqrt{d_o}$ and $\sigma_t(\bar{W}_{\rm WN}) = 0$ for $t=2,\dots,T$, which means that  $\bar{W}_{\rm WN}$ is of rank one. Similar thing can be said to the Frobenius normalization (See the appendix for more details). Using such $\bar{W}_{\rm WN}$ corresponds to using only one feature to discriminate the model probability distribution from the target. In order to retain as much norm of the input as possible and hence to make the discriminator more sensitive, one would hope to make the norm of $\bar{W}_{\rm WN}{\bm h}$ large. For weight normalization, however, this comes at the cost of reducing the rank and hence the number of features to be used for the discriminator. Thus, there is a conflict of interests between weight normalization and our desire to use as many features as possible to distinguish the generator distribution from the target distribution. The former interest often reigns over the other in many cases, inadvertently diminishing the number of features to be used by the discriminators. Consequently, the algorithm would produce a rather arbitrary model distribution that matches the target distribution only at select few features. Weight clipping~\cite[]{arjovsky2017wasserstein} also suffers from same pitfall.

Our spectral normalization, on the other hand, do not suffer from such a conflict in interest. Note that the Lipschitz constant of a linear operator is determined only by the maximum singular value. In other words, the spectral norm is independent of rank. Thus, unlike the weight normalization, our spectral normalization allows the parameter matrix to use as many features as possible while satisfying local $1$-Lipschitz constraint. Our spectral normalization leaves more freedom in choosing the number of singular components (features) to feed to the next layer of the discriminator.

\cite{brock2016neural} introduced orthonormal regularization on each weight to stabilize the training of GANs. In their work, \cite{brock2016neural} augmented the adversarial objective function by adding the following term:

\begin{align}
	\|W^{\rm T} W-I\|_F^2.
\end{align}

While this seems to serve the same purpose as spectral normalization, orthonormal regularization are mathematically quite different from our spectral normalization because the orthonormal regularization destroys the information about the spectrum by setting all the singular values to one. On the other hand, spectral normalization only scales the spectrum so that the its maximum will be one. 

\cite{gulrajani2017improved} used Gradient penalty method in combination with WGAN. In their work, they placed $K$-Lipschitz constant on the discriminator by augmenting the objective function with the regularizer that rewards the function for having local 1-Lipschitz constant (i.e. $\|\nabla_{\hat{{\bm x}}}f\|_2=1$)  at discrete sets of points of the form $\hat{{\bm x}} := \epsilon \tilde {\bm x} + (1- \epsilon) {\bm x}$ generated by interpolating a sample $\tilde {\bm x}$ from generative distribution and a sample ${\bm x}$ from the data distribution. While this rather straightforward approach does not suffer from the problems we mentioned above regarding the effective dimension of the feature space, the approach has an obvious weakness of being heavily dependent on the support of the current generative distribution. As a matter of course, the generative distribution and its support gradually changes in the course of the training, and this can destabilize the effect of such regularization. In fact, we empirically observed that a high learning rate can destabilize the performance of WGAN-GP. On the contrary, our spectral normalization regularizes the function the operator space, and the effect of the regularization is more stable with respect to the choice of the batch. Training with our spectral normalization does not easily destabilize with aggressive learning rate. Moreover, WGAN-GP requires more computational cost than our spectral normalization with single-step power iteration, because the computation of  $\|\nabla_{\hat{{\bm x}}}f\|_2$ requires one whole round of forward and backward propagation. In the appendix section, we compare the computational cost of the two methods for the same number of updates.

\section{Experiments}
In order to evaluate the efficacy of our approach and investigate the reason behind its efficacy, we conducted a set of extensive experiments of unsupervised image generation on CIFAR-10~\cite[]{torralba200880} and STL-10~\cite[]{coates2011analysis}, and compared our method against other normalization techniques. To see how our method fares against large dataset, we also applied our method on ILSVRC2012 dataset (ImageNet)~\cite[]{ILSVRC15} as well. This section is structured as follows. First, we will discuss the objective functions we used to train the architecture, and then we will describe the optimization settings we used in the experiments. We will then explain two performance measures on the images to evaluate the images produced by the trained generators. Finally, we will summarize our results on CIFAR-10, STL-10, and ImageNet.

As for the architecture of the discriminator and generator, we used convolutional neural networks. Also, for the evaluation of the spectral norm for the convolutional weight $W \in \mathbb{R}^{d_{\rm out}\times d_{\rm in} \times h \times w}$, we treated the operator as a 2-D matrix of dimension $d_{\rm out}\times (d_{\rm in} h w)$\footnote{Note that, since we are conducting the convolution discretely, the spectral norm will depend on the size of the stride and padding. However, the answer will only differ by some predefined $K$.}. We trained the parameters of the generator with batch normalization~\cite[]{ioffe2015batch}.
We refer the readers to Table~\ref{tab:cnn_models} in the appendix section for more details of the architectures.

For all methods other than WGAN-GP, we used the following standard objective function for the adversarial loss:

\begin{align}
	V(G, D) := \mathop{\mathrm{E}}_{x\sim q_{\rm data}({\bm x})} [ \log D({\bm x})] +  \mathop{\mathrm{E}}_{{\bm z}\sim p({\bm z})} [\log(1-D(G({\bm z})))], \label{eq:advloss}
\end{align}

where ${\bm z} \in \mathbb{R}^{d_z}$ is a latent variable, $p({\bm z})$ is the standard normal distribution $\mathcal{N}(0, I)$, and $G: \mathbb{R}^{d_z} \to \mathbb{R}^{d_0}$ is a deterministic generator function. We set $d_z$ to 128 for all of our experiments. For the updates of $G$, we used the alternate cost proposed by~\cite{goodfellow2014generative}
$-\mathop{\mathrm{E}}_{{\bm z}\sim p({\bm z})} [\log(D(G({\bm z})))]$ as used in~\cite{goodfellow2014generative} and~\cite{warde2016improving}. For the updates of $D$, we used the original cost defined in~\eqref{eq:advloss}.
We also tested the performance of the algorithm with the so-called hinge loss, which is given by

\begin{align}
	 V_D(\hat{G}, D) &= \mathop{\mathrm{E}}_{{\bm x}\sim q_{\rm data}({\bm x})}\left[{\rm min}\left(0, -1+D({\bm x})\right)\right] + \mathop{\mathrm{E}}_{{\bm z}\sim p({\bm z})} \left[{\rm min}\left(0, -1-D\left(\hat{G}({\bm z})\right)\right)\right]\\
     V_G(G, \hat{D}) &= -\mathop{\mathrm{E}}_{{\bm z}\sim p({\bm z})}\left[\hat{D}\left(G({\bm z})\right)\right], \label{eq:hinge}
\end{align}

respectively for the discriminator and the generator. Optimizing these objectives is equivalent to minimizing the so-called reverse KL divergence : ${\rm KL}[p_g||q_{\rm data}]$. This type of loss has been already proposed and used in~\cite{lim2017geometric, tran2017deep}. The algorithm based on the hinge loss also showed good performance when evaluated with inception score and FID. For Wasserstein GANs with gradient penalty (WGAN-GP)~\cite[]{gulrajani2017improved}, we used the following objective function: $V(G, D) := \mathop{\mathrm{E}}_{{\bm x} \sim q_{\rm data}} [ D({\bm x})] -  \mathop{\mathrm{E}}_{{\bm z}\sim p({\bm z})} [D(G({\bm z}))] - \lambda \mathop{\mathrm{E}}_{\hat{{\bm x}} \sim p_{\hat{{\bm x}}}}[(\|\nabla_{\hat{{\bm x}}}D(\hat{{\bm x}})\|_2-1)^2] \label{eq:wgan_gp}$,
where the regularization term is the one we introduced in the appendix section~\ref{subsubsec:wgan-gp}.

For quantitative assessment of generated examples, we used \textit{inception score}~\cite[]{salimans2016improved} and \textit{Fr\'echet inception distance} (FID)~\cite[]{heusel2017gans}.
Please see Appendix~\ref{apdsubsec:performance_measures} for the details of each score.

\subsection{Results on CIFAR10 and STL-10} In this section, we report the accuracy of the spectral normalization (we use the abbreviation: SN-GAN for the spectrally normalized GANs) during the training, and the dependence of the algorithm's performance on the hyperparmeters of the optimizer. We also compare the performance quality of the algorithm against those of other regularization/normalization techniques for the discriminator networks, including: Weight clipping~\cite[]{arjovsky2017wasserstein}, WGAN-GP~\cite[]{gulrajani2017improved}, batch-normalization (BN)~\cite[]{ioffe2015batch}, layer normalization (LN)~\cite[]{ba2016layer}, weight normalization (WN)~\cite[]{salimans2016weight} and orthonormal regularization (\textit{orthonormal})~\cite[]{brock2016neural}.
In order to evaluate the stand-alone efficacy of the gradient penalty, we also applied the gradient penalty term to the standard adversarial loss of GANs~\eqref{eq:advloss}.
We would refer to this method as `GAN-GP'. For weight clipping, we followed the original work \cite{arjovsky2017wasserstein} and set the clipping constant $c$ at 0.01 for the convolutional weight of each layer. For gradient penalty, we set $\lambda$ to 10, as suggested in~\cite{gulrajani2017improved}. For \textit{orthonormal}, we initialized the each weight of $D$ with a randomly selected orthonormal operator and trained GANs with the objective function augmented with the regularization term used in \cite{brock2016neural}. For all comparative studies throughout, we excluded the multiplier parameter ${\bm \gamma}$ in the weight normalization method, as well as in batch normalization and layer normalization method. This was done in order to prevent the methods from overtly violating the Lipschitz condition. When we experimented \textit{with} different multiplier parameter, we were in fact not able to achieve any improvement.

For optimization, we used the Adam optimizer~\cite{kingma2014adam} in all of our experiments. We tested with 6 settings for (1) $n_{\rm dis}$, the number of updates of the discriminator per one update of the generator and (2) learning rate $\alpha$ and the first and second order momentum parameters ($\beta_1$, $\beta_2$) of Adam. We list the details of these settings in Table~\ref{tab:settings} in the appendix section. Out of these 6 settings, A, B, and C are the settings used in previous representative works. The purpose of the settings D, E, and F is to the evaluate the performance of the algorithms implemented with more aggressive learning rates. For the details of the architectures of convolutional networks deployed in the generator and the discriminator, we refer the readers to Table~\ref{tab:cnn_models} in the appendix section. The number of updates for GAN generator were 100K for all experiments, unless otherwise noted.

Firstly, we inspected the spectral norm of each layer during the training to make sure that our spectral normalization procedure is indeed serving its purpose. As we can see in the Figure \ref{fig:snorm} in the \ref{apdsubsec:acc_sn}, the spectral norms of these layers floats around $1$--$1.05$ region throughout the training. Please see Appendix \ref{apdsubsec:acc_sn} for more details.

In Figures~\ref{fig:incpscores} and ~\ref{fig:fids} we show the inception scores of each method with the settings A--F. We can see that spectral normalization is relatively robust with aggressive learning rates and momentum parameters. WGAN-GP fails to train good GANs at high learning rates and high momentum parameters on both CIFAR-10 and STL-10. Orthonormal regularization performed poorly for the setting E on the STL-10, but performed slightly better than our method with the optimal setting. These results suggests that our method is more robust than other methods with respect to the change in the setting of the training. Also, the optimal performance of weight normalization was inferior to both WGAN-GP and spectral normalization on STL-10, which consists of more diverse examples than CIFAR-10. Best scores of spectral normalization are better than almost all other methods on both CIFAR-10 and STL-10.

In Tables~\ref{tab:incepscores}, we show the inception scores of the different methods with optimal settings on CIFAR-10 and STL-10 dataset. We see that SN-GANs performed better than almost all contemporaries on the optimal settings. SN-GANs performed even better with hinge loss~\eqref{eq:hinge}.\footnote{As for STL-10, we also ran SN-GANs over twice time longer iterations because it did not seem to converge. Yet still, this elongated training sequence still completes before WGAN-GP with original iteration size because the optimal setting of SN-GANs (setting B, $n_{dis} =1$) is computationally light.}. For the training with same number of iterations, SN-GANs fell behind orthonormal regularization for STL-10. For more detailed comparison between orthonormal regularization and spectral normalization, please see section \ref{sec:compare_agstORTH}. 

In Figure~\ref{fig:images} we show the images produced by the generators trained with WGAN-GP, weight normalization, and spectral normalization. SN-GANs were consistently better than GANs with weight normalization in terms of the quality of generated images. To be more precise, as we mentioned in Section~\ref{sec:snvsothers}, the set of images generated by spectral normalization was clearer and more diverse than the images produced by the weight normalization. We can also see that WGAN-GP failed to train good GANs with high learning rates and high momentums (D,E and F). The generated images with GAN-GP, batch normalization, and layer normalization is shown in Figure~\ref{fig:cifar10images_others} in the appendix section.

We also compared our algorithm against multiple benchmark methods ans summarized the results on the bottom half of the Table~\ref{tab:incepscores}.
We also tested the performance of our method on ResNet based GANs used in \cite{gulrajani2017improved}. Please note that all methods listed thereof are all different in both optimization methods and the architecture of the model. Please see Table~\ref{tab:resnets_cifar10} and ~\ref{tab:resnets_stl} in the appendix section for the detail network architectures. Our implementation of our algorithm was able to perform better than almost all the predecessors in the performance.  

\footnotetext{For our ResNet experiments, we trained the same architecture with multiple random seeds for weight initialization and produced models with different parameters. We then generated 5000 images 10 times and computed the average inception score for each model. The values for ResNet on the table are the mean and standard deviation of the score computed over the set of models trained with different seeds.}

\subsubsection{Analysis of SN-GANs}

\paragraph{Singular values analysis on the weights of the discriminator $D$} In Figure~\ref{fig:svhist}, we show the squared singular values of the weight matrices in the final discriminator $D$ produced by each method using the parameter that yielded the best inception score. As we predicted in Section~\ref{sec:snvsothers}, the singular values of the first to fifth layers trained with weight clipping and weight normalization concentrate on a few components. That is, the weight matrices of these layers tend to be rank deficit. On the other hand, the singular values of the weight matrices in those layers trained with spectral normalization is more broadly distributed. When the goal is to distinguish a pair of probability distributions on the low-dimensional nonlinear data manifold embedded in a high dimensional space, rank deficiencies in lower layers can be especially fatal. Outputs of lower layers have gone through only a few sets of rectified linear transformations, which means that they tend to lie on the space that is linear in most parts. Marginalizing out many features of the input distribution in such space can result in oversimplified discriminator. We can actually confirm the effect of this phenomenon on the generated images especially in Figure~\ref{fig:stlimages}.
The images generated with spectral normalization is more diverse and complex than those generated with weight normalization.

\paragraph{Training time} On CIFAR-10, SN-GANs is slightly slower than weight normalization (about 110 $\sim$ 120\% computational time), but significantly faster than WGAN-GP. As we mentioned in Section~\ref{sec:snvsothers}, WGAN-GP is slower than other methods because WGAN-GP needs to calculate the gradient of gradient norm $\|\nabla_{{\bm x}} D\|_2$.
For STL-10, the computational time of SN-GANs is almost the same as vanilla GANs, because the relative computational cost of the power iteration~\eqref{eq:poweriter} is negligible when compared to the cost of forward and backward propagation on CIFAR-10 (images size of STL-10 is larger ($48 \times 48$)). Please see Figure~\ref{fig:cptime} in the appendix section for the actual computational time.

\subsubsection{Comparison between SN-GANs and orthonormal regularization} \label{sec:compare_agstORTH} 

In order to highlight the difference between our spectral normalization and orthonormal regularization, we conducted an additional set of experiments. As we explained in Section~\ref{sec:snvsothers}, orthonormal regularization is different from our method in that it destroys the spectral information and puts equal emphasis on all feature dimensions, including the ones that 'shall' be weeded out in the training process. To see the extent of its possibly detrimental effect, we experimented by increasing the dimension of the feature space \footnote{More precisely, we simply increased the input dimension and the output dimension by the same factor. In Figure ~\ref{fig:varying_fmaps}, `relative size' = 1.0 implies that the layer structure is the same as the original.}, especially at the final layer (7th conv) for which the training with our spectral normalization prefers relatively small feature space (dimension $<100$; see Figure~\ref{fig:svhist_stl}). 
As for the setting of the training, we selected the parameters for which the orthonormal regularization performed optimally. The figure~\ref{fig:varying_fmaps} shows the result of our experiments. As we predicted, the performance of the orthonormal regularization deteriorates as we increase the dimension of the feature maps at the final layer. Our SN-GANs, on the other hand, does not falter with this modification of the architecture. Thus, at least in this perspective, we may such that our method is more robust with respect to the change of the network architecture.

\subsection{Image Generation on ImageNet}

To show that our method remains effective on a large high dimensional dataset, we also applied our method to the training of conditional GANs on ILRSVRC2012 dataset with 1000 classes, each consisting of approximately 1300 images, which we compressed to $128 \times 128$ pixels. Regarding the adversarial loss for conditional GANs, we used practically the same formulation used in \cite{mirza2014conditional}, except that we replaced the standard GANs loss with hinge loss \eqref{eq:hinge}.
Please see Appendix~\ref{apdsubsec:exp_imagenet} for the details of experimental settings.

GANs without normalization and GANs with layer normalization collapsed in the beginning of training and failed to produce any meaningful images. GANs with orthonormal normalization \cite{brock2016neural} and our spectral normalization, on the other hand, was able to produce images. The inception score of the orthonormal normalization however plateaued around 20Kth iterations, while SN kept improving even afterward (Figure~\ref{fig:imagenet_incp}.) 
To our knowledge, our research is the first of its kind in succeeding to produce decent images from ImageNet dataset with a \textit{single} pair of a discriminator and a generator (Figure~\ref{fig:generated_images_imagenet}).
To measure the degree of mode-collapse, we followed the footstep of~\cite{odena2016conditional} and computed the intra MS-SSIM~\cite{odena2016conditional} for pairs of independently generated GANs images of each class. We see that our SN-GANs ((intra MS-SSIM)=0.101) is suffering less from the mode-collapse than AC-GANs ((intra MS-SSIM)$\sim$0.25).

To ensure that the superiority of our method is not limited within our specific setting, we also compared the performance of SN-GANs against \textit{orthonormal regularization} on conditional GANs with \textit{projection discriminator}~\cite[]{miyato2018cgans} as well as the standard (unconditional) GANs. In our experiments, SN-GANs achieved better performance than \textit{orthonormal regularization} for the both settings (See Figure~\ref{fig:imagenet_incp_others} in the appendix section).

\section{Conclusion}
This paper proposes spectral normalization as a stabilizer of training of GANs. When we apply spectral normalization to the GANs on image generation tasks, the generated examples are more diverse than the conventional weight normalization and achieve better or comparative inception scores relative to previous studies. The method imposes global regularization on the discriminator as opposed to local regularization introduced by WGAN-GP, and can possibly used in combinations. In the future work, we would like to further investigate where our methods stand amongst other methods on more theoretical basis, and experiment our algorithm on larger and more complex datasets.

\subsubsection*{Acknowledgments}
 We would like to thank the members of Preferred Networks, Inc., particularly Shin-ichi Maeda, Eiichi Matsumoto, Masaki Watanabe and Keisuke Yahata for insightful comments and discussions. We also would like to thank anonymous reviewers and commenters on the OpenReview forum for insightful discussions.

\end{document}
